\section{Introduction}
\label{sec:introduction}

Drug-target interaction (DTI) prediction is a fundamental problem in computational drug discovery because it connects candidate molecules with potential therapeutic targets before expensive biochemical validation. Reliable DTI models can accelerate chemical-genomics screening, target identification, and drug repositioning by narrowing the search space of plausible compound-protein pairs \citep{keiser2007relating,yamanishi2008prediction,liu2007bindingdb}. The task is challenging because interaction signals depend jointly on molecular chemistry, three-dimensional ligand geometry, protein sequence context, and the pair-specific compatibility between the two modalities. As a result, an effective model should not only encode drugs and targets accurately in isolation, but also preserve the interaction structure that determines whether a given pair is likely to bind.

Existing DTI methods have progressively expanded the representation space used for this task. Early network-inference and similarity-based methods integrated chemical and genomic spaces to infer unknown drug-target links \citep{yamanishi2008prediction}. Sequence-based deep models such as DeepDTA use neural encoders for SMILES strings and protein sequences, while graph-based models such as GraphDTA encode molecular topology with graph neural networks \citep{ozturk2018deepdta,nguyen2021graphdta}. Transformer-style and attention-based models further improve compound-protein modeling by learning sequence, substructure, or bilinear interaction patterns \citep{chen2020transformercpi,huang2021moltrans,bai2023drugban,peng2024bindti}. More recent multimodal frameworks combine pretrained molecular and protein language models, graph encoders, equivariant geometric encoders, and cross-modal fusion modules to exploit complementary drug and target views \citep{chithrananda2020chemberta,elnaggar2021prottrans,satorras2021egnn,ji2026ldmdti}. These developments show that DTI prediction benefits from jointly using sequence semantics, molecular topology, geometric information, and explicit interaction modeling.

Despite this progress, two limitations remain underexplored. First, existing multimodal DTI models usually encode molecular geometry as a drug-side representation before cross-modal fusion \citep{satorras2021egnn,ji2026ldmdti}; however, the same ligand may expose different relevant conformer cues when paired with different proteins. This weakens the ability of the model to select conformer information under the actual target context. Second, long-range protein sequence modeling and drug-target matching are often entangled in a single fusion design. Although structured state-space models provide efficient long-sequence modeling \citep{gu2022s4,gu2023mamba}, directly applying sequence refinement to mixed drug-protein tokens can blur modality structure, while replacing explicit cross-modal attention may remove the pairwise matching bias that DTI prediction requires. These issues suggest that target-conditioned geometry and protein sequence refinement should be introduced carefully, without disrupting the final interaction module.

To address these limitations, we propose TAMR-DTI, a target-aware multimodal framework built on a strong BiIntention-based DTI backbone. TAMR-DTI introduces bidirectional drug-target modulation: protein context guides conformer-level drug aggregation, and fused ligand information conditions protein token encoding through a lightweight modulation path. In addition, we insert a gated Mamba residual block only on the ordered protein token stream before BiIntention \citep{gu2023mamba}. This design uses state-space modeling where biological token order is meaningful, while preserving BiIntention as the explicit drug-target matching component. In this way, TAMR-DTI strengthens target-aware representation learning and long-range protein refinement without imposing artificial sequential dynamics on heterogeneous drug tokens.

The main contributions of this work are summarized as follows:
\begin{itemize}
  \item We formulate a target-aware modulation strategy that couples conformer aggregation with protein context and conditions protein token encoding on ligand information.
  \item We introduce a protein-side Mamba residual refinement module that improves ordered protein representations while retaining explicit BiIntention-based drug-target matching.
  \item We evaluate TAMR-DTI on BindingDB, BioSNAP, Human, and C. elegans under repeated-seed experimental protocols, showing consistent ranking-metric improvements across all four benchmarks.
  \item We conduct comparative and ablation studies to verify the effectiveness of target-aware modulation and protein-only state-space refinement within a multimodal DTI framework.
\end{itemize}
